\section{Scope and Genericity}
\label{sec:scope}

This section records what the theorem does and does not claim.

The 6-vertex example analyzed in Section \ref{sec:example} is a small, reproducible check of the construction: the generator masks, the labeled violation matrix, the ranks, and the final count can all be inspected directly.

The generic mirror separation assumption (Definition \ref{def:mirror-sep}) is essential for the exactness of the rank count. It excludes accidental geometric coincidences. In non-generic configurations, the actual number of solutions can be larger than the generic rank count. This happens when an apparently violating edge is preserved accidentally, for example because a point lies unexpectedly on a reflection mirror. Such cases form a proper algebraic exceptional set of measure zero in the coordinate parameter space. Analogous almost-always qualifications appear in existing DDGP and DMDGP counting results \cite{liberti2014number,abud2024impossible}.

The case $K=1$ requires separate treatment because the algebraic bracket-separation arguments used to establish the generic presentation lemma for $K \ge 2$ degenerate on the line.

\subsection{Unification of the Combinatorial DDGP}
\label{sec:combinatorial-ddgp}

We now check that the rank-count formula agrees with a known positive counting result. 

Abud et al. (2024) introduced the subclass of \emph{Combinatorial DDGP}, defined as those DDGP instances where every predecessor set $U_j$ induces a clique of size $K$ in $G$ \cite{abud2024impossible}. They proved that, for this subclass and in the absence of pruning edges, solution counting is combinatorial and yields exactly $2^{n-K}$ incongruent solutions. This is Corollary 3.6 of Abud et al. and also generalizes the DMDGP no-pruning case.

The rank-count formula includes this result as a direct special case:
\begin{enumerate}
    \item When every predecessor set $U_j$ induces a clique, any active discretization edge $\{a, b\} \in E_D[L_P]$ between predecessors of a vertex has both of its endpoints within the clique.
    \item Consequently, if a generator support crosses the boundary, the fixed endpoint is guaranteed to lie on the mirror clique. Thus, no active edges can be violated under any branch transformations.
    \item The labeled violation matrix $V$ is therefore identically the zero matrix ($V = 0$).
    \item Substituting $V = 0$ into our rank formula:
    \[
    \dim \mathcal{K}_F
    =
    \operatorname{rank}\begin{bmatrix}M\\ 0\end{bmatrix} - \operatorname{rank}(0)
    =
    \operatorname{rank}(M) - 0
    =
    n - K,
    \]
    where $M$ has full rank $n-K$ because the cone generators form a triangular system of masks.
\end{enumerate}
This yields exactly $|\Xi_F| = 2^{n-K}$, recovering the positive result of Abud et al. (2024). Thus the rank-count formula agrees with clique-based DMDGP and Combinatorial DDGP counts while also applying to local non-clique DDGP subproblems satisfying the theorem's hypotheses.
